\scanpage{122}
\mspage{6}
% archive sheet number: 121
\uline{Corollaire 2} Considérons un carré \deleted{$h$-cartésien}
\diagram{\begin{tikzcd}[ampersand replacement=\&,execute at end picture={
\node[anchor=east] at ([xshift=-5pt]current bounding box.west) {$\deleted{X}\unsure{Y|X}=\Biggl($};
\node[anchor=west] at ([xshift=5pt]current bounding box.east) {$\Biggr)$};
}]
\deleted{X_1}Y_1 \arrow[d] \& \deleted{X_3}Y_3 \arrow[l] \arrow[d] \\
\deleted{X_0}Y_0 \& \deleted{X_2}Y_2 \arrow[l]
\end{tikzcd}}
,
et un objet $T$ de $\unsure{\mathscr H|\mathcal H|H}$, d'où (posant $\operatorname{Hom}(T,X)=\deleted{X}Y(T)$) un carré comm. des $(\mathrm{Ens})$
\diagram{\begin{tikzcd}[ampersand replacement=\&]
\deleted{X_1}Y_1(T) \arrow[d] \& \deleted{X_3}Y_3(T) \arrow[l] \arrow[d] \\
\deleted{X_0}Y_0(T) \& \deleted{X_2}Y_2(T) \arrow[l]
\end{tikzcd}}
, et par suite
\[
\deleted{X_3}Y_3(T)\longrightarrow \deleted{X_1}Y_1(T)\mathbin{\times}_{\deleted{X_0}Y_0(T)}\deleted{X_2}Y_2(T).\tag*{(*)}
\]
Si $\deleted{X}Y$ est $h$-cartésien, cette application \inserted{$(*)$} est \uline{surjective} quel que soit $T$. \deleted{(\unsure{mais|ceci} \unsure{que}}

On applique $\operatorname{Cart}(\unsure{6|5})$ à $Y$ et au carré constant \inserted{$X$} de valeur $T$
\diagram{\begin{tikzcd}[ampersand replacement=\&,execute at end picture={
\node[anchor=east] at ([xshift=-5pt]current bounding box.west) {$X=$};
}]
T \arrow[d,"\mathrm{id}"'] \& T \arrow[l,"\mathrm{id}"'] \arrow[d,"\mathrm{id}"] \\
T \& T \arrow[l,"\mathrm{id}"]
\end{tikzcd}}
, en notant que la donnée de $p=(p_0,p_1,p_2)$ de
\diagram{\begin{tikzcd}[ampersand replacement=\&]
T \arrow[d] \& \\
T \& T \arrow[l]
\end{tikzcd}}
dans
\diagram{\begin{tikzcd}[ampersand replacement=\&]
Y_1 \arrow[d] \& \\
Y_0 \& Y_2 \arrow[l]
\end{tikzcd}}
revient à la donnée d'un élément du deuxième mb de $(*)$, et que \deleted{\unsure{l'existence d'un}}
\diagram{\begin{tikzcd}[ampersand replacement=\&]
Y_1 \arrow[d] \& T \arrow[l,"p_1"'] \arrow[d,"p_2"] \arrow[dl,dashed,"p_0"'] \\
\unsure{Y_0|\deleted{Y_0}} \& \unsure{Y_2|\deleted{Y_2}} \arrow[l]
\end{tikzcd}}
éléments \deleted{$p_3$} de $Y_3(T)$ qui s'envoient (par $*$) dans cet élément ne sont autres que les $p_3$ dont l'axiome $\operatorname{Cart}(6)$ affirme l'existence.

\uline{NB} On notera que l'application \inserted{surjective} $(*)$ n'est en général pas bijective, i.e. le carré $h$-cartésien n'est en général pas cartésien.
